\documentclass[12pt]{article}

%Paquetes a utilizarse
\usepackage[width=7in, height=9.5in, top=0.75in, papersize={8.5in,11in}]{geometry}
\usepackage[spanish]{babel} 
\decimalpoint
\usepackage[utf8]{inputenc}
\usepackage{bbding}
\usepackage[colorlinks = true, linkcolor = blue, urlcolor = BlueViolet, citecolor = OliveGreen]{hyperref}
\usepackage{graphicx}
\usepackage{amssymb,amsthm,amsmath}
\usepackage{enumerate}
\usepackage{array,multicol,multirow}
\usepackage{xcolor}
\usepackage{fancybox,tcolorbox}
\usepackage{caption,subcaption,float,tabularx}
\usepackage{enumitem}

\theoremstyle{definition}
\newtheorem{corolario}{Corolario}
\newtheorem{lema}[corolario]{Lema}
\newtheorem{proposicion}[corolario]{Proposición}
\newtheorem{teorema}[corolario]{Teorema}
\newtheorem{propiedad}[corolario]{Propiedad}
\newtheorem*{observacion}{Observación}
\newtheorem{definicion}{Definición}
\newtheorem*{demostracion}{Demostración}
\newtheorem{ejemplo}{Ejemplo}
\newtheorem{problema}{Problema}
\newtheorem*{solucion}{Solución}
\newtheorem{ejercicio}{\PencilRightDown \  Ejercicio}
\newtheorem{step}{Paso}
\newtheorem{credito}{Crédito}

\usepackage{tikz}
\usetikzlibrary{arrows.meta,babel,calc,positioning}

\renewcommand{\arraystretch}{1.5}
\providecommand{\abs}[1]{\lvert#1\rvert}
\providecommand{\norm}[1]{\lVert#1\rVert}

\renewcommand{\tabularxcolumn}[1]{m{#1}}
\newcommand{\Evaluacion}[4]{
\setcounter{ejercicio}{0}
\noindent\begin{tabular}{lcr}
	\includegraphics[height=3cm]{Logos/logo-UES.png}\hspace{2.5em}
	&
	\includegraphics[height=2.75cm]{Logos/logo-PJT.png}
	& 
	\hspace{2.5em}\includegraphics[height=2.75cm]{Logos/logo-MINEDUCYT.png}
\end{tabular}

\hfill

\begin{center}
    
    UNIVERSIDAD DE EL SALVADOR
    \\PROGRAMA JÓVENES TALENTO
    \\FDTC 2022
    \\#2
    \\Nivel Olímpico C de Matemáticas

\end{center}

\begin{center}
    #1
\end{center}

%\textbf{Nombre}: \enspace\hrulefill

#3

\input{#4}
\newpage
}

\newtheorem{obs}{Observación}

%\usepackage[margin=2.5cm]{geometry}
%\usepackage{wasysym}
%\usepackage{stmaryrd,textcomp}
%\usepackage{pgf,tikz}
%\usetikzlibrary{arrows}

\parskip = 2mm   %%%% genera un espacio de X mm entre lo párrafos
\parindent = 3mm
\usepackage{multicol}
\usepackage{iwona}

\newcommand{\tema}{Semana 1}
\newcommand{\fecha}{Sábado, 9 de Diciembre de 2023}
\newcommand{\sesion}{Semanal 1}

\begin{document}
%\thispagestyle{empty}
%\newpage
\thispagestyle{empty}

\begin{figure}[h] 
	\begin{minipage}[b]{0.26\textwidth}
		\begin{center}
			\includegraphics[height=3cm]{Logos/UES.png}
			\par\end{center}
	\end{minipage} 
	\begin{minipage}[b]{0.46\textwidth}
		\begin{center}
			UNIVERSIDAD DE EL SALVADOR\\ [0.1cm]
			PROGRAMA JÓVENES TALENTO\\ [0.1cm]
	        FDTC 2022\\ [0.1cm]
                NIVEL 5\\ [0.1cm]
			COMBINATORIA 
			\par\end{center}
	\end{minipage} 
	\begin{minipage}[b]{0.05\textwidth}
		\begin{center}
			\includegraphics[height=2cm]{Logos/LOGO PJT.png}
			\par\end{center}
	\end{minipage}
\end{figure}

\begin{center}
    \begin{tabular}{p{4.5cm} p{7cm} p{4.5cm}}
        \tema & \centering\fecha & \hfill\sesion
    \end{tabular}
\end{center}
Nombre: \rule{15cm}{1pt}\newline

\textbf{INDICACIONES:} Resuelve cada problema de manera clara y ordenada, dejando constancia de tus procedimientos.
\vspace{0.2cm}
\begin{enumerate}
    \item (25\%) Demuestre que, dados $5$ puntos en el plano con coordenadas enteras, siempre hay dos tales que su punto medio tiene coordenadas enteras. Tenga en cuenta que las coordenadas del punto medio de dos puntos, $A=(x_1,y_1)$ y $B=(x_2,y_2)$, están dadas por $\displaystyle M=\left(\frac{x_1 + x_2}{2}, \frac{y_1 + y_2}{2}\right)$.
    \item (25\%) Para los enteros positivos  $1,2,3,\dots ,n$ hay 11660 desarreglos donde $1,2,3,4,5$ aparecen en los primeros cinco lugares. ¿Cuál es el valor de $n$?
    \item (25\%) Imagina que tienes cuatro monedas diferentes: una de 1 euro, una de 50 centavos, una de 20 centavos y una de 10 centavos. ¿De cuántas maneras diferentes puedes colocar estas monedas en una fila si la moneda de 50 centavos siempre debe estar separada de la moneda de 1 euro por al menos una moneda?
    \item (25\%) Utilizando el método de contar de dos formas, demuestra la siguiente propiedad del Hockey Stick (o Christmas Stocking Theorem) en el Triángulo de Pascal: 

\textit{La suma de los números combinatorios a lo largo de una diagonal del Triángulo de Pascal, comenzando con \( \displaystyle\binom{r}{r} \) y extendiéndose hacia abajo hasta \( \displaystyle\binom{n}{r} \), es igual a \( \displaystyle \binom{n+1}{r+1} \).}

Matemáticamente, se expresa como:
\[ \binom{r}{r} + \binom{r+1}{r} + \binom{r+2}{r} + \cdots + \binom{n}{r} = \binom{n+1}{r+1} \]
\end{enumerate}

\textbf{Crédito Extra:} Dada la expresión \[\left(7x^2+\frac{1}{2x}\right)^{25}\]encuentre la cantidad de términos de su expansión, la suma de los coeficientes de su expansión, el coeficiente de $x^{13}$ si es posible, y encuentre el i-ésimo término

\end{document}